% 分类目录：dl
\documentclass[12pt, a4paper]{article}
\usepackage[margin=2.5cm]{geometry}
\usepackage{ctex}
\usepackage{amsmath, amssymb}
\usepackage{listings}
\usepackage{xcolor}
\usepackage{tabularx}
\usepackage{hyperref}

\lstset{
    language=Python,
    basicstyle=\ttfamily\small,
    keywordstyle=\color{blue},
    commentstyle=\color{green!50!black},
    stringstyle=\color{red},
    showstringspaces=false,
    breaklines=true,
    numbers=left,
    numberstyle=\tiny\color{gray},
    frame=single
}

\title{知识点：特征工程 (Feature Engineering)}
\author{}
\date{}

\begin{document}
\maketitle

\section{核心定义}
\textbf{中文定义}：特征工程是利用领域知识对原始数据进行处理，以创建能够使机器学习/深度学习模型更好地执行的特征的过程。它包含特征的选择、提取、构造和变换。在深度学习时代，虽然神经网络具有自动提取特征的能力，但在处理结构化数据（Tabular Data）和特定领域任务时，高质量的手工特征依然至关重要。

\textbf{英文定义}：Feature engineering is the process of using domain knowledge to process raw data to create features that make machine learning or deep learning models perform better. It involves feature selection, extraction, construction, and transformation. Although deep neural networks can extract features automatically, high-quality handcrafted features remain essential for structured data and domain-specific tasks.

\section{特征选择 (Feature Selection)}

\begin{itemize}
    \item \textbf{过滤式 (Filter)}：直接评估特征与标签的相关性。
        \begin{itemize}
            \item \textbf{皮尔逊相关系数} (Pearson Correlation)：衡量线性相关性。
            \item \textbf{卡方检验} (Chi-Square)：衡量分类变量的相关性。
        \end{itemize}
    \item \textbf{包裹式 (Wrapper)}：将模型视为黑盒，通过循环迭代筛选最优组合，如 \textbf{RFE (递归特征消除)}。
    \item \textbf{嵌入式 (Embedded)}：将特征选择集成在训练过程中。
        \begin{itemize}
            \item \textbf{L1 正则化 (Lasso)}：使不重要特征的权重变为 0，从而实现自动选择。
            \item \textbf{树模型重要性}：利用随机森林或 XGBoost 的特征重要度。
        \end{itemize}
\end{itemize}

\section{特征构造 (Feature Construction)}
\begin{itemize}
    \item \textbf{组合特征} (Interaction Terms)：将两个原始特征相乘或相除，例如“点击次数/曝光次数 = 点击率”。
    \item \textbf{时序特征}：滞后项（Lag）、滑动窗口统计（Rolling Mean）、日期分解（星期、小时等）。
    \item \textbf{聚合特征}：结合分组统计，如“用户过去 30 天的平均消费额”。
\end{itemize}

\section{特征提取 (Feature Extraction)}
通过数学变换将高维空间映射到低维：
\begin{itemize}
    \item \textbf{PCA (主成分分析)}：无监督，最大化投影方差。
    \item \textbf{LDA (线性判别分析)}：有监督，最大化类间方差且最小化类内方差。
    \item \textbf{Autoencoders (自编码器)}：利用神经网络进行非线性的特征压缩。
\end{itemize}

\section{特征变换 (Feature Transformation)}
\begin{enumerate}
    \item \textbf{偏态分布优化}：使用 \textbf{Log 变换} 或 \textbf{Box-Cox 变换} 使重尾分布趋于正态。
    \item \textbf{高基数类别处理}：
        \begin{itemize}
            \item \textbf{均值编码 (Mean Encoding)}：用对应标签的均值替换类别，需警惕过拟合。
            \item \textbf{哈希技巧 (Hashing Trick)}：固定维度，适用于特征空间巨大的文本处理。
        \end{itemize}
\end{enumerate}

\section{深度学习中的特征工程}
\textbf{问题}：既然神经网络能自动提取特征，特征工程是否过时了？
\\ \textbf{答}：在视觉和语音领域，CNN 和 RNN 确实消除了大部分传统特征。但在以下场景，特征工程依然是核心：
\begin{itemize}
    \item \textbf{结构化数据}：对于包含大量离散、长尾字段的表格数据，显式地构造交叉特征远比加深神经网络更有效。
    \item \textbf{领域约束}：将物理定律或先验经验以特征形式输入，可以显著降低模型对样本量的需求。
    \item \textbf{Embedding}：深度学习本身就是一种强力的特征工程，它将离散符号转化为连续语义特征。
\end{itemize}

\section{关键代码片段 (Python)}
\begin{lstlisting}
from sklearn.feature_selection import SelectFromModel
from sklearn.linear_model import Lasso
import pandas as pd

# 1. 嵌入式特征选择 (L1)
lasso = Lasso(alpha=0.1)
selector = SelectFromModel(lasso)
selector.fit(X, y)
selected_features = X.columns[selector.get_support()]

# 2. 构造多项式特征
from sklearn.preprocessing import PolynomialFeatures
poly = PolynomialFeatures(degree=2, interaction_only=True)
X_poly = poly.fit_transform(X)
\end{lstlisting}

\section{深度面试题}

\textbf{问：什么是特征穿越 (Feature Leakage)？}\\
答：指特征中包含了本不应知道的标签信息。例如，在预测是否流失的任务中，包含了“流失日期”或“账户关闭记录”作为特征。这会导致训练集表现完美，但上线即失效。

\textbf{问：多重共线性 (Multicollinearity) 对模型有何影响？}\\
答：在基于梯度的优化中，它会导致 Hessian 矩阵接近奇异，使模型系数极其不稳定，解释性变差。可以通过 VIF (方差膨胀因子) 检验，并剔除 VIF > 10 的冗余特征。

\section{参考文献}
\begin{enumerate}
    \item Kuhn, M., \& Johnson, K. (2013). \textit{Applied Predictive Modeling}. Springer.
    \item Guyon, I., \& Elisseeff, A. (2003). \textit{An Introduction to Variable and Feature Selection}. Journal of Machine Learning Research.
    \item Dong, G., \& Liu, H. (2018). \textit{Feature Engineering for Machine Learning and Data Analytics}. CRC Press.
\end{enumerate}

\end{document}
